\section{Conclusion}\label{section7}


In conclusion, this work provides a comprehensive theoretical foundation for understanding the capabilities of multi-round reasoning in auto-regressive large language models. We have systematically explored the approximation, learnability, and generalization properties of these models, highlighting their potential to solve complex problems through iterative reasoning. 

\textbf{Findings.} We demonstrate that Transformers with a finite context window can serve as universal approximators for Turing-computable functions, offering insights into their robust problem-solving capabilities in real-world tasks. Additionally, we have extended the PAC learning framework to account for sequence generation tasks, revealing the complexities involved in learning long sequences when context exceeds the model's window size. 
Moreover, our analysis of the generalization error in multi-round generation reveals that, without proper interventions, error accumulation could lead to divergence in the model's outputs. Techniques like Chain-of-Thought offer viable strategies to mitigate this, ensuring that the generated sequences remain within expected bounds. 

\textbf{Practical Insights.} Our contributions not only advance our theoretical understanding of auto-regressive generative language models but also provide practical insights into improving model performance through multi-round reasoning interventions. 
 For the model training stage, in order to reduce the sample complexity of training on long sequences, one can consider providing some decomposition methods for very long and complex task sequences during language model training, so that the long sequences are decomposed into multiple rounds of training on short sequences. In the inference process of the model, when we design a method that makes the model perform multi-round thinking, we should give more consideration to how to interrupt the propagation of cumulative errors to make the generated content more in line with our expectations.